%%%%%%%%%%%%%%%%%%%%%%%%%%%%%%%%%%%%%%%%%%%%%%%%%%%%%%%%%%%%%%%%%%%%%%%%%%%%%%%%%%
\begin{frame}[fragile]\frametitle{}
\begin{center}
{\Large Introduction}
\end{center}
\end{frame}


%%%%%%%%%%%%%%%%%%%%%%%%%%%%%%%%%%%%%%%%%%%%%%%%%%%%%%%%%%%
\begin{frame}[fragile]\frametitle{Pace of AI Innovation}

\begin{center}
\includegraphics[width=\linewidth,keepaspectratio]{promptengg57}

{\tiny (Ref: The Complete Prompt Engineering for AI Bootcamp (2023))}
\end{center}

\textbf{Key insight:} Notice the Stable Diffusion line (blue), it is nearly vertical. AI capabilities are not growing gradually; they are exploding. Prompt Engineering is the skill that lets you \emph{harness} this growth today.
				
\end{frame}

%%%%%%%%%%%%%%%%%%%%%%%%%%%%%%%%%%%%%%%%%%%%%%%%%%%%%%%%%%%
\begin{frame}[fragile]\frametitle{Prompt Engineering Is Becoming a Universal Skill}

\begin{center}
\includegraphics[width=\linewidth,keepaspectratio]{promptengg58}

{\tiny (Ref: The Complete Prompt Engineering for AI Bootcamp (2023))}
\end{center}		

\textbf{Just like Google Search} became a skill everyone needed in the 2000s, Prompt Engineering is the essential skill of this decade, for professionals, students, and everyday users alike.
		
\end{frame}


%%%%%%%%%%%%%%%%%%%%%%%%%%%%%%%%%%%%%%%%%%%%%%%%%%%%%%%%%%%
\begin{frame}[fragile]\frametitle{What is a Prompt?}

\begin{columns}
    \begin{column}[T]{0.6\linewidth}
		\begin{center}
		\includegraphics[width=\linewidth,keepaspectratio]{promptengg55}

		{\tiny (Ref: The Complete Prompt Engineering for AI Bootcamp (2023))}
		\end{center}	
    \end{column}
    \begin{column}[T]{0.4\linewidth}
    	\textbf{In plain language:} A prompt is simply the text you type when talking to an AI.
    	
    	\medskip
    	Think of it like giving instructions to a very capable but very literal assistant.
    	
    	\medskip
		\textbf{Intuitive:} A cue given to the AI so it understands what you want.
		
		\medskip
		\textbf{Technical:} Additional text added to the input that helps the model make better use of its knowledge.
    \end{column}
  \end{columns}
\end{frame}


%%%%%%%%%%%%%%%%%%%%%%%%%%%%%%%%%%%%%%%%%%%%%%%%%%%%%%%%%%%
\begin{frame}[fragile]\frametitle{Is an LLM Just Glorified Auto-complete?}

\begin{center}
\includegraphics[width=0.85\linewidth,keepaspectratio]{promptengg42}

{\tiny (Ref: The Fourth Paradigm of Modern Natural Language Processing, Pengfei Liu)}
\end{center}

\medskip
\textbf{Yes, but at a massive scale.} Just like your phone suggests the next word when you type, an LLM predicts the next word, except it was trained on billions of sentences, so its ``suggestions'' are often impressively intelligent and context-aware.

\end{frame}


%%%%%%%%%%%%%%%%%%%%%%%%%%%%%%%%%%%%%%%%%%%%%%%%%%%%%%%%%%%
\begin{frame}[fragile]\frametitle{What is Prompt Engineering?}

\begin{columns}
    \begin{column}[T]{0.6\linewidth}
		\begin{center}
		\includegraphics[width=\linewidth,keepaspectratio]{promptengg56}

		{\tiny (Ref: The Complete Prompt Engineering for AI Bootcamp (2023))}
		\end{center}	
    \end{column}
    \begin{column}[T]{0.4\linewidth}
		\textbf{Prompt Engineering} is the art and science of crafting your input text so that the AI reliably gives you the output you actually want.

		\medskip
		Think of it like learning how to ask the right question at the right time, same as you would with a very knowledgeable expert.
    \end{column}
  \end{columns}
\end{frame}


%%%%%%%%%%%%%%%%%%%%%%%%%%%%%%%%%%%%%%%%%%%%%%%%%%%%%%%%%%%
\begin{frame}[fragile]\frametitle{Why Does Prompt Engineering Matter?}

% \begin{columns}
    % \begin{column}[T]{0.55\linewidth}
    \begin{itemize}
    \item \textbf{Same AI, very different results:} The exact same AI model can give an excellent answer or a useless one, the only difference is \emph{how} you asked.
    \item \textbf{No coding needed:} You do not need to be a programmer to get great results from AI. A well-crafted prompt is all you need.
    \item \textbf{Unlocks advanced AI capabilities:} Enables complex reasoning, code generation, content creation, and decision support.
    \item \textbf{Reduces hallucinations:} A poorly worded prompt invites confident but wrong answers. Good prompts anchor the AI to facts.
    \item \textbf{Critical for research \& business:} Helps test AI limits and build reliable applications.
    \end{itemize}
    {\tiny (Ref: Prompt Engineering, A lecture by DAIR.AI)}
    % \end{column}
    % \begin{column}[T]{0.45\linewidth}
		% \begin{center}
		% \includegraphics[width=\linewidth,keepaspectratio]{promptengg2}
		% {\tiny (Ref: Cohere https://docs.cohere.ai/docs/prompt-engineering)}
		% \end{center}
    % \end{column}
% \end{columns}

\end{frame}



%%%%%%%%%%%%%%%%%%%%%%%%%%%%%%%%%%%%%%%%%%%%%%%%%%%%%%%%%%%
\begin{frame}[fragile]\frametitle{Why Word Choice Matters: A Motivating Example}

\begin{itemize}
\item Ask the AI: \lstinline|What is 1,000,000 * 9,000?|
\item GPT-3 sometimes answers: \textbf{9,000,000} \quad (\emph{wrong!})

\medskip
\item Now ask: \lstinline|What is 1,000,000 * 9,000? Make sure to put the right amount of zeros, even if there are many:|
\item GPT-3 answers: \textbf{9,000,000,000} \quad (\emph{correct!})

\medskip
\item \textbf{Why does this happen?} The AI is not a calculator — it is a language model. It ``reads'' your words and generates a plausible-sounding response. Adding the hint about zeros nudged it to be more careful.

\medskip
\item \textbf{This is the whole point of Prompt Engineering} — figuring out \emph{how} to word your request so the AI does exactly what you want.
\end{itemize}

{\tiny (Ref: https://learnprompting.org/docs/basics/prompting)}
\end{frame}


%%%%%%%%%%%%%%%%%%%%%%%%%%%%%%%%%%%%%%%%%%%%%%%%%%%%%%%%%%%
\begin{frame}[fragile]\frametitle{What Can Good Prompting Achieve?}

\begin{columns}
    \begin{column}[T]{0.55\linewidth}
    \begin{itemize}
    \item \textbf{Avoid Inconsistent Results:} Without a good prompt, the same AI might give you a different answer every time, sometimes correct, sometimes wrong. Good prompts make outputs predictable and reliable.
    
    \medskip
    \item \textbf{Use Advanced Techniques:} Combine styles (``paint this like Picasso but with Van Gogh's colours''), create personalised avatars, or build creative content at scale.
    
    \medskip
    \item \textbf{Scale Creative Output:} Generate dozens of variations of an image, article, or code snippet, programmatically, by changing just a few words in your prompt.
    \end{itemize}
    \end{column}
    \begin{column}[T]{0.45\linewidth}
		\begin{center}
		\includegraphics[width=\linewidth,keepaspectratio]{promptengg59}
		{\tiny Avoiding hallucination: notice image with missing arms?}

		\medskip
		\includegraphics[width=\linewidth,keepaspectratio]{promptengg60}
		{\tiny Mixing artistic styles with prompting.}
		
		{\tiny (Ref: The Complete Prompt Engineering for AI Bootcamp (2023))}

		\end{center}
    \end{column}
\end{columns}


\end{frame}



%%%%%%%%%%%%%%%%%%%%%%%%%%%%%%%%%%%%%%%%%%%%%%%%%%%%%%%%%%%%%%%%%%%%%%%%%%%%%%%%%%
\begin{frame}[fragile]\frametitle{}
\begin{center}
{\Large Applications}
\end{center}
\end{frame}

%%%%%%%%%%%%%%%%%%%%%%%%%%%%%%%%%%%%%%%%%%%%%%%%%%%%%%%%%%%
\begin{frame}[fragile]\frametitle{Text Summarization}
\begin{itemize}
    \item \textbf{Task:} Condense a long piece of text into a short, clear summary.
    \item \textbf{How:} Paste the text and tell the AI how short the summary should be.
\end{itemize}
\begin{lstlisting}[language=HTML]
Prompt:
Antibiotics are a type of medication used to treat bacterial 
infections. They work by either killing the bacteria or 
preventing them from reproducing, allowing the body's immune 
system to fight off the infection. They are not effective 
against viral infections, and using them inappropriately 
can lead to antibiotic resistance.
Explain the above in one sentence:

Output:
Antibiotics treat bacterial infections by killing or 
stopping bacteria from reproducing, but don't work on 
viruses, and overuse causes resistance.
\end{lstlisting}
\end{frame}

%%%%%%%%%%%%%%%%%%%%%%%%%%%%%%%%%%%%%%%%%%%%%%%%%%%%%%%%%%%
\begin{frame}[fragile]\frametitle{Information Extraction}
\begin{itemize}
    \item \textbf{Task:} Pull out a specific piece of information from a larger block of text — like a search inside a document.
    \item \textbf{How:} Give the AI the text and tell it exactly what to look for.
\end{itemize}
\begin{lstlisting}[language=HTML]
Prompt:
Author-contribution statements in research papers should 
state clearly whether authors used AI technologies such as 
ChatGPT in the preparation of their manuscript...

Mention the large language model based product mentioned 
in the paragraph above:

Output:
The large language model based product mentioned is ChatGPT.
\end{lstlisting}
\end{frame}

%%%%%%%%%%%%%%%%%%%%%%%%%%%%%%%%%%%%%%%%%%%%%%%%%%%%%%%%%%%
\begin{frame}[fragile]\frametitle{Question Answering}
\begin{itemize}
    \item \textbf{Task:} Answer a specific question using only the context provided — the AI does not guess; it reads what you gave it.
    \item \textbf{How:} Give context + question. Ask it to say ``Unsure'' if the answer is not in the context.
\end{itemize}
\begin{lstlisting}[language=HTML]
Answer the question based on the context below. 
Keep the answer short. Respond "Unsure about answer" 
if not sure about the answer.

Context: Teplizumab traces its roots to a New Jersey drug 
company called Ortho Pharmaceutical. There, scientists 
generated an early version of the antibody, dubbed OKT3.

Question: What was OKT3 originally sourced from?

Output:
Mice.
\end{lstlisting}
\end{frame}

%%%%%%%%%%%%%%%%%%%%%%%%%%%%%%%%%%%%%%%%%%%%%%%%%%%%%%%%%%%
\begin{frame}[fragile]\frametitle{Code Generation}
\begin{itemize}
    \item \textbf{Task:} Generate working code just from a plain-English description — no programming knowledge needed to describe \emph{what} you want.
\end{itemize}
\begin{lstlisting}[language=HTML]
Prompt:
/*
Ask the user for their name and say "Hello"
*/
Output:
let name = prompt("What is your name?");
console.log(`Hello, ${name}!`);

--------------------------------------------------
Prompt:
Table departments, columns = [DepartmentId, DepartmentName]
Table students, columns = [DepartmentId, StudentId, StudentName]
Create a MySQL query for all students in Computer Science

Output:
SELECT StudentId, StudentName FROM students 
WHERE DepartmentId IN 
  (SELECT DepartmentId FROM departments 
   WHERE DepartmentName = 'Computer Science');
\end{lstlisting}
\end{frame}

%%%%%%%%%%%%%%%%%%%%%%%%%%%%%%%%%%%%%%%%%%%%%%%%%%%%%%%%%%%
% IMPROVED: Merged the two Reasoning slides, added plain-English explanation of the failure
\begin{frame}[fragile]\frametitle{Reasoning — and Where It Can Go Wrong}
\begin{itemize}
    \item LLMs can do arithmetic but they sometimes get it wrong because they are \textbf{predicting likely text}, not computing.
    \item \textbf{Fix:} Break the problem into steps in your prompt.
\end{itemize}
\begin{lstlisting}[language=HTML]
Prompt (naive):
What is 9,000 * 9,000?
Output: 81,000,000   [correct here, but unreliable on harder sums]

--------------------------------------------------
Prompt (step-by-step):
The odd numbers in this group add up to an even number: 
15, 32, 5, 13, 82, 7, 1. 
Solve by breaking the problem into steps. First, identify 
the odd numbers, add them, and indicate whether the result 
is odd or even.

Output:
Odd numbers: 15, 5, 13, 7, 1
Sum: 41 
41 is an odd number.
\end{lstlisting}
\end{frame}

%%%%%%%%%%%%%%%%%%%%%%%%%%%%%%%%%%%%%%%%%%%%%%%%%%%%%%%%%%%
\begin{frame}[fragile]\frametitle{Image Generation}

\begin{center}
\includegraphics[width=\linewidth,keepaspectratio]{promptengg15}

{\tiny (Ref: Prompt Engineering Sudalai Rajkumar)}
\end{center}		
		
\textbf{Models / Tools:} DALL-E, Midjourney, Stable Diffusion.

\medskip
You describe the image in words — the AI generates it. The more precise your description (style, lighting, composition, subject), the closer the result is to what you imagined.

\end{frame}


%%%%%%%%%%%%%%%%%%%%%%%%%%%%%%%%%%%%%%%%%%%%%%%%%%%%%%%%%%%
\begin{frame}[fragile]\frametitle{Calling a Prompt Programmatically (Python)}

\begin{lstlisting}[language=Python]
# Requires: pip install openai python-dotenv
from openai import OpenAI
import os
from dotenv import load_dotenv, find_dotenv

load_dotenv(find_dotenv())
client = OpenAI(api_key=os.getenv('OPENAI_API_KEY'))

def get_completion(prompt, model="gpt-4o"):
    response = client.chat.completions.create(
        model=model,
        messages=[{"role": "user", "content": prompt}],
        temperature=0,  # 0 = deterministic, 1 = creative
    )
    return response.choices[0].message.content

# Example usage
result = get_completion("Explain photosynthesis in one sentence.")
print(result)
\end{lstlisting}

% {\tiny Note: Uses the modern \texttt{openai} SDK (v1.0+). The legacy \texttt{ChatCompletion.create()} API is deprecated.}
		
\end{frame}


%%%%%%%%%%%%%%%%%%%%%%%%%%%%%%%%%%%%%%%%%%%%%%%%%%%%%%%%%%%
\begin{frame}[fragile]\frametitle{Key Settings That Control AI Output}

When calling an AI via an API or playground, you can tune several ``dials'':

\begin{itemize}
\item \textbf{Temperature}: Controls \emph{creativity vs. consistency}.
  \begin{itemize}
  \item \texttt{0} = Always picks the most likely word. Same answer every time. Great for facts.
  \item \texttt{1} = Picks more surprising words. Different answer each time. Great for creative writing.
  \item Think of it like a volume knob for randomness.
  \end{itemize}
\item \textbf{Top-P (Nucleus Sampling)}: Controls \emph{how wide a vocabulary} the AI considers.
  \begin{itemize}
  \item \texttt{Top-P = 0.9} means: only pick from words that together cover 90\% of the likely choices.
  \item Lower values make output more focused; higher values more varied.
  \item \textbf{Use either Temperature or Top-P, not both} (OpenAI recommendation).
  \end{itemize}
\item \textbf{Max Tokens}: The maximum length of the response (1 token $\approx$ 0.75 words).
\item \textbf{Frequency / Presence Penalty}: Reduces repetition in the output.
\end{itemize}

\end{frame}


%%%%%%%%%%%%%%%%%%%%%%%%%%%%%%%%%%%%%%%%%%%%%%%%%%%%%%%%%%%
\begin{frame}[fragile]\frametitle{Advanced Prompting Techniques: Overview}

\begin{center}
\includegraphics[width=0.75\linewidth,keepaspectratio]{promptengg97}
\end{center}				

{\tiny (Ref: Applied LLMs Mastery 2024 -- Aishwarya Reganti)}

\end{frame}


